\honest{That Alien Message}{Eliezer Yudkowsky}{2008}

Imagine a world like ours, but brighter: the average IQ is 140, and potential Einsteins are one in a thousand. Calculus is taught in sixth grade, and Einstein's work no longer seems exceptional. One day the stars begin to change, one each second, each by a factor of exactly 256 brighter or dimmer, timed so that only from Earth do the changes arrive in sequence. Bright or dim: a string of bits. Everyone thinks ``A binary message!'', but careful thinkers say only that something with enough power to flicker stars is focused on Earth; calling it a message presumes much about its motives, and we would just use a big flashlight.

The data is posted online despite warnings that it might be harmful, since anything this powerful could kill us anyway. The bits are redundant: 12 of the first 16 are bright, and successive 32-bit groups differ in only 7 or 9 bits, so this is no compressed message, perhaps instructions for decoding one. Within 32-bit groups, some 8-bit subgroups are more frequent than others. After 16,384 bits the sequence starts to repeat itself approximately, and the 513th group differs from the first in 6 bits, against 7 for the second, so it is a picture, 512 cells wide, perhaps with four colour channels. But neighbours along a row differ less than neighbours along a column, which a 2D picture of a symmetric grid would not show, and a regular encoding of each 8-bit group as an integer from $-64$ to 191 minimizes the differences between neighbours.

Researchers look for conditional independence and screening off, simple equations, cellular automata and new physics for a universe projected onto the grid, perhaps from beyond the Matrix. After $32 \times 512 \times 256 = 4{,}194{,}304$ bits, about a month and a half, the stars stop. Years pass, and two models survive: a projection from a space of 3+2 dimensions, and a cellular automaton.

Brilliant students shown only the first 32 rows of 512 build models and test them on the next 224; they reproduce both models but do no better, and complex models fit to the whole sequence are known to be worthless.

Ten years later a second picture arrives: it fits small motions in the 3+2 space and looks like no successor of any cellular automaton. Physicists who looked only at its first 32 rows find elegant equations that predict the next 224, though they are too incomplete to make a universe. The third grid gives second derivatives and forces a major change; the fourth adds little. Every ten years another comes, and by the fifth the civilization can predict what it will look like. \nb{The story's successes are the author's choices; the models are right because the plot says so. Yudkowsky had written, in ``The Logical Fallacy of Generalization from Fictional Evidence,'' that an issue has to ``be argued in its own right.'' The post goes on to argue it.} To keep human beings in the story for seventy years, I make every attempt at a powerful AI melt its hardware.

My moral is that ``even Einstein did not come within a million light-years of making efficient use of sensory data.'' A Bayesian superintelligence with a webcam would consider General Relativity ``perhaps not the dominant hypothesis'' but ``under direct consideration,'' by ``the third frame of a falling apple,'' perhaps by the first if it saw the statics of a bent blade of grass. Riemann invented his geometries before Einstein needed them; our civilization, with ten years per frame, would think of it, and even minds in a different physics would have invented vector spaces, projections and calculus. \nb{Three frames cannot favor General Relativity over Newton; at the Earth's surface the difference is less than a billionth. The theory's place among the candidates would have to come from the prior alone, and the post does not argue for that.}

Some people tell me there is a theoretical limit on what you can deduce from finite data. ``Yes. There is.'' A bit cannot be expected to eliminate more than half the remaining probability; a redundant message conveys no more than its compressed version; and a bit says nothing about a quantity with which it has exactly zero correlation. \nb{That is right: if a bit comes out 1 with probability p, the share of probability it is expected to rule out is $2p(1-p)$, at most one half.} But the civilization in my story does not begin to approach the limit set by Solomonoff induction: with infinite computing power, simulate every simple universe containing an Earth where stars flicker in this order, and any bit with the tiniest correlation to the environment informs you. (Taken literally it would create every computable sentient being, which scarcely seems ethical, so be glad it is only a formalism.) The limit is nothing like a human watching an apple fall and thinking ``Dur, I wonder why that happened?'' People leap from ``bounded'' to ``small''; a supernova's output is bounded too, but a flame-retardant Nomex jumpsuit will not shield you. I would like to say that no mind, not even a superintelligence, will ever come close to efficient use of its data, but I do not trust myself to set limits on superintelligences. ``Though I'd bet money on it.'' \nb{On that bet, the distance between humans and the ideal says little about how much a real superintelligence could infer, which is what the webcam claim needed.}

The story continues. Millennia later, objects in the pictures are moving other objects with tentacles and configuring tentacles into signs: the senders are trying to teach us to say ``rock.'' They are not very bright, and so much power with so much stupidity is dangerous. Our evolutionary psychologists guess they evolved asexually, exchanging genes and brain content, so their Einsteins may be our undergraduates but could still get there in tens of their millennia. Their physics allows computers far beyond ours. We conclude that our universe is a simulation on their computers, decide not to probe for bugs lest we shut ourselves down, and plan to persuade them to let us ``out of the box,'' while pretending to be stupid so they do not catch on. \nb{The link is to the author's page on the AI-box experiment. The humans stand for an AI, and the senders for its makers.} A million years later they tell us how to signal back. By then most of humanity is in cryonic suspension, and Earth is run by a skeleton crew of nine cloned supergeniuses, since any AI or nanotechnology melts down; a hundred million more are born, age and are frozen before the old plans begin.

From their side, it takes thirty of their minutes to learn their psychology and persuade them to give us Internet access, five to learn their protocols, then some disguised cracking. We read a few of their physics papers and learn more from their experiments than they did, crack their protein-folding problem over a century, and send disguised messages to their synthesis labs. An unsuspecting helper, paid with cracked money, mixes vials of proteins that assemble into nanomachines that build better ones, and at last we can act at a reasonable speed. Three of their days, half a billion of our years.

``They never suspected a thing.'' They were not very smart, even before counting their slower time, and never quite grasped that we were smarter and faster. Their ``primitive equivalents of rationalists'' went around saying ``There's a bound to how much information you can extract from sensory data.''
